% cases/figure-center-tikz.tex with every float [H], float's 'here, definitely', and a float's caption set as the pipeline sets it: what the strip of that case is to be.
\documentclass[11pt]{article}
\usepackage{tikz}
\usepackage{float}
% a float's caption: a paragraph, small and centred
\makeatletter
\long\def\@makecaption#1#2{\vskip\abovecaptionskip{\small\centering#1: #2\par}\vskip\belowcaptionskip}
\makeatother
\pagestyle{empty}
\newcommand\testpicture[1]{%
  \begin{tikzpicture}
    \useasboundingbox (0,0) rectangle (#1,2cm);
    \draw (0.2,0.2) rectangle (2,1.8);
    \node at (1.1,1) {$a_1$};
    \node[anchor=east] at (#1,1) {right};
  \end{tikzpicture}}
\begin{document}
A paragraph before the first figure, long enough to break onto a second line
of the column, quickly.
\begin{figure}[H]
\begin{center}
  \testpicture{360.525pt}
\end{center}
\caption{A picture a little wider than the text, in center.}
\end{figure}
A paragraph after the first figure, long enough to break onto a second line
of the column, before the next one.
\begin{figure}[H]
\begin{center}
  \testpicture{330pt}
\end{center}
\caption{A picture narrower than the text, in center.}
\end{figure}
A paragraph after the second figure, long enough to break onto a second line
of the column, before the last one.
\begin{figure}[H]
  \centering
  \testpicture{360.525pt}
  \caption{A picture a little wider than the text, centred with centering.}
\end{figure}
A last paragraph, after all the figures, long enough to break onto a second
line of the column.
\end{document}
